    \section{Dimension}
        \subsection{Definition of dimension}
            \hskip \parindent
            \hskip \parindent
            \par Let $V$ be a finitely generated vector space (f.d.v.s.) over a field $\mathbb{F}$, and let $B=\{v_1,v_2,\cdots,v_n\}$ be a basis for $V$. The number $n$ is called the dimension of $V$.
            \begin{equation}
                dim_{\mathbb{F}}V=n
            \end{equation}
        \subsection{Fundamental lemma of linear algebra}
            \hskip \parindent
            \par Let $\{v_1,v_2,\cdots,v_n\}$ be a linearly independent set of vectors in a vector space $V$ over a field $\mathbb{F}$. Then any subset of $\operatorname{span}\{v_1,v_2,\cdots,v_n\}$ with more than $n$ vectors is linearly dependent.
        \subsection{Theorem: elements of basis have same number\label{sec:theorem_dimension}}
            \hskip \parindent
            \par If $\{v_1,v_2,\cdots,v_m\}$ and $\{w_1,w_2,\cdots,w_n\}$ are basis for a vector space $V$ over a field $\mathbb{F}$, then $n=m$.
            \par Proof see section "Proofs".
        \subsection{Example of dimension}
            \begin{itemize}
                \item $dim_{\mathbb{R}}\mathbb{R}^n=n$
                \item $dim_{\mathbb{F}}M_{m\times n}(\mathbb{F})=mn$
                \item $dim_{\mathbb{R}}\mathbb{C}^1=2$
                \item $dim_{\mathbb{R}}\mathbb{C}^n=2n$
            \end{itemize}
        \subsection{Definition of column space and row space}
            \hskip \parindent
            \par $A \in M_{m\times n}(\mathbb{F})$, $A=[v_1|v_2|\cdots|v_n]$ or $A =\begin{bmatrix} w_1\\ w_2\\ \cdots\\ w_m\end{bmatrix}$ where $v_i \in \mathbb{F}^m$, $1 \leq i \leq n$ and $w_j \in \mathbb{F}^n$, $1 \leq j \leq m$.
            \par Column-space = span$\{v_1,v_2,\cdots,v_n\}$ $\subseteq \mathbb{F}^m$.
            \par Row-space = span$\{w_1,w_2,\cdots,w_m\}$ $\subseteq \mathbb{F}^n$.
        \subsection{Definition of rank}
            \hskip \parindent
            \par Let $A$ be in $M_{m\times n}(\mathbb{F})$. We difine the rank of $A$ as the dimension of the subspace of $\mathbb{F}^m$ generated by the columns (column space) of $A$. We denote this number by rk $A$ or rank $A$.
        \subsection{Remark: rank and non-zero rows/pivot columns}
            \hskip \parindent
            \par The definition of the rank is equivalent to the number of non-zero rows in the row-reduced echelon form $R$ that is row-equivalent to $A$ or the number of pivot columns in $R$.
        \subsection{Two properties}
            \subsubsection{Row equivalent matrices have the same rank}
                \hskip \parindent
                \par If $A$ and $A'$ are row equivalent, then they have the same rank.
            \subsubsection{rank $A$ = rank $R$}
                \hskip \parindent
                \par If $R$ is a row-reduced echelon form that is row equivalent to $A$, then rank $A$ = rank $R$.
        \subsection{Theorem: dimension of subspace is less than or equal to the dimension of the space \label{sec:theorem_dimension_of_subspace}}
            \hskip \parindent
            \par Let $W$ be a subspace of a f.d.v.s. $V$ over a field $\mathbb{F}$. Then $W$ is also a f.d.v.s. and $dim_{\mathbb{F}}W \leq dim_{\mathbb{F}}V$.
            \par The proof see section "Proofs".